% ============================================================================
% LISTE COMPLÈTE DES FIGURES ET TABLEAUX — RAPPORT GUIDNI (Ahmed Addali)
% Structurée selon le modèle du rapport d'Amal Rahmeni
% ============================================================================

\documentclass[12pt,a4paper]{article}
\usepackage[utf8]{inputenc}
\usepackage[T1]{fontenc}
\usepackage[french]{babel}
\usepackage{geometry}
\usepackage{booktabs}
\usepackage{longtable}
\usepackage{array}
\usepackage{xcolor}

\geometry{margin=2.5cm}

\title{Liste des Figures et Tableaux du Rapport Guidni\\
\large Système de Planification de Voyage Intelligent\\
\small Structuré selon le modèle d'Amal Rahmeni}
\author{Ahmed Addali}
\date{\today}

\begin{document}

\maketitle

\tableofcontents

\newpage

% ============================================================================
% CHAPITRE 1 — CONTEXTE GÉNÉRAL
% ============================================================================
\section{Chapitre 1 : Contexte Général}

\subsection*{Figures attendues}

\begin{longtable}{p{2cm} p{12cm}}
    \textbf{Référence} & \textbf{Description} \\
    \midrule
    \endhead
    
    Figure 1.1 & Logo de l'entreprise/organisme d'accueil (ISITCOM) \\
    Figure 1.2 & Logo de la plateforme Guidni \\
    Figure 1.3 & Les services de la plateforme Guidni (vue d'ensemble) \\
    Figure 1.4 & Types d'agents LLM (ReAct simple, LangGraph, Agent avec RAG) \\
    Figure 1.5 & Exemple d'agent ReAct simple (architecture basique) \\
    Figure 1.6 & Exemple d'agent LangGraph (graphe d'états à 4 nœuds) \\
    Figure 1.7 & Exemple d'agent avec RAG (pipeline de retrieval) \\
    Figure 1.8 & Types de méthodologies agiles (2TUP / Scrum / Kanban) \\
    Figure 1.9 & Processus 2TUP appliqué à Guidni \\
    Figure 1.10 & Processus détaillé de 2TUP (phases et livrables) \\
    Figure 1.11 & Notre équipe projet (photos des membres si disponible) \\
    \bottomrule
\end{longtable}

\subsection*{Tableaux attendus}

\begin{longtable}{p{2cm} p{12cm}}
    \textbf{Référence} & \textbf{Description} \\
    \midrule
    \endhead
    
    Tableau 1.1 & Identité de l'organisme d'accueil (ISITCOM : raison sociale, forme juridique, adresse, téléphone, site web) \\
    Tableau 1.2 & Tableau des tâches (fonctionnalités à concevoir et à implémenter pour Guidni) \\
    Tableau 1.3 & Tableau comparatif entre les méthodologies (2TUP vs Scrum vs Kanban avec colonnes : Description / Points Forts / Points Faibles) \\
    \bottomrule
\end{longtable}

\subsection*{Statut de création}

\begin{itemize}
    \item[\textcolor{green}{\checkmark}] Figure 1.1 : Logo ISITCOM — \textbf{À CRÉER} (remplacer par logo officiel)
    \item[\textcolor{green}{\checkmark}] Figure 1.2 : Logo Guidni — \textbf{À CRÉER} (design à réaliser)
    \item[\textcolor{orange}{!}] Figure 1.3 : Services Guidni — \textbf{PARTIELLEMENT EXISTANT} (adapter screen\_01\_djerba\_tourism.png)
    \item[\textcolor{green}{\checkmark}] Figure 1.4 : Types d'agents LLM — \textbf{À CRÉER} (diagramme comparatif)
    \item[\textcolor{green}{\checkmark}] Figure 1.5 : Agent ReAct simple — \textbf{À CRÉER} (diagramme UML)
    \item[\textcolor{green}{\checkmark}] Figure 1.6 : Agent LangGraph — \textbf{EXISTANT} (fig\_02\_langgraph\_graph.png)
    \item[\textcolor{green}{\checkmark}] Figure 1.7 : Agent avec RAG — \textbf{EXISTANT} (fig\_02\_rag\_architecture.png)
    \item[\textcolor{green}{\checkmark}] Figure 1.8 : Méthodologies agiles — \textbf{EXISTANT} (fig\_01\_2tup\_methodology.png à adapter)
    \item[\textcolor{green}{\checkmark}] Figure 1.9 : Processus 2TUP — \textbf{EXISTANT} (fig\_01\_2tup\_methodology.png)
    \item[\textcolor{green}{\checkmark}] Figure 1.10 : Processus détaillé 2TUP — \textbf{À ADAPTER} (à partir de fig\_01\_2tup\_methodology.png)
    \item[\textcolor{red}{\times}] Figure 1.11 : Équipe projet — \textbf{OPTIONNEL} (photos non disponibles)
    
    \item[\textcolor{green}{\checkmark}] Tableau 1.1 : Identité ISITCOM — \textbf{À CRÉER} (informations institutionnelles)
    \item[\textcolor{green}{\checkmark}] Tableau 1.2 : Tâches Guidni — \textbf{À CRÉER} (liste des fonctionnalités)
    \item[\textcolor{green}{\checkmark}] Tableau 1.3 : Comparatif méthodologies — \textbf{À CRÉER} (2TUP/Scrum/Kanban)
\end{itemize}


% ============================================================================
% CHAPITRE 2 — LANCEMENT DE PROJET
% ============================================================================
\section{Chapitre 2 : Lancement de Projet}

\subsection*{Figures attendues}

\begin{longtable}{p{2cm} p{12cm}}
    \textbf{Référence} & \textbf{Description} \\
    \midrule
    \endhead
    
    Figure 2.1 & Diagramme de contexte (acteurs et système Guidni) \\
    Figure 2.2 & Architecture du pipeline RAG/LLM (flux de traitement) \\
    Figure 2.3 & Domaine des LLMs et agents (cartographie conceptuelle) \\
    Figure 2.4 & Architecture du Transformer (décodeur) \\
    Figure 2.5 & Architecture LangGraph (graphe d'états avec 4 nœuds) \\
    Figure 2.6 & Comparaison ReAct vs LangGraph (avantages/inconvénients) \\
    Figure 2.7 & Architecture fonctionnelle de Guidni \\
    Figure 2.8 & Architecture logique générale de Guidni \\
    Figure 2.9 & Diagramme de Gantt / Organigramme Guidni \\
    Figure 2.10 & Diagramme de cas d'utilisation global \\
    Figure 2.11 & Definition of Done (DoD) pour Guidni \\
    Figure 2.12 & Logo Python + description (5-6 lignes) \\
    Figure 2.13 & Logo FastAPI + description (5-6 lignes) \\
    Figure 2.14 & Logo PostgreSQL + description (5-6 lignes) \\
    Figure 2.15 & Logo React/Next.js + description (5-6 lignes) \\
    Figure 2.16 & Logo LangGraph + description (5-6 lignes) \\
    Figure 2.17 & Logo OpenAI API + description (5-6 lignes) \\
    Figure 2.18 & Logo FAISS/pgvector + description (5-6 lignes) \\
    Figure 2.19 & Logo SQLAlchemy + description (5-6 lignes) \\
    Figure 2.20 & Logo APScheduler + description (5-6 lignes) \\
    \bottomrule
\end{longtable}

\subsection*{Tableaux attendus}

\begin{longtable}{p{2cm} p{12cm}}
    \textbf{Référence} & \textbf{Description} \\
    \midrule
    \endhead
    
    Tableau 2.1 & Étude comparative entre modèles (LLM vs RAG vs LinUCB vs approches classiques) \\
    Tableau 2.2 & Product Backlog (ID / User Story / Priorité MoSCoW) \\
    Tableau 2.3 & Planification des différents Sprints (Release / Sprint / Objectif / Période) \\
    Tableau 2.4 & Caractéristiques de l'ordinateur de développement \\
    \bottomrule
\end{longtable}

\subsection*{Statut de création}

\begin{itemize}
    \item[\textcolor{green}{\checkmark}] Figure 2.1 : Diagramme de contexte — \textbf{EXISTANT} (fig\_03\_use\_case\_global.png peut être adapté)
    \item[\textcolor{green}{\checkmark}] Figure 2.2 : Pipeline RAG/LLM — \textbf{EXISTANT} (fig\_02\_rag\_architecture.png)
    \item[\textcolor{green}{\checkmark}] Figure 2.3 : Domaine LLMs — \textbf{À CRÉER} (cartographie mentale)
    \item[\textcolor{green}{\checkmark}] Figure 2.4 : Architecture Transformer — \textbf{EXISTANT} (fig\_02\_transformer\_attention.png)
    \item[\textcolor{green}{\checkmark}] Figure 2.5 : Architecture LangGraph — \textbf{EXISTANT} (fig\_02\_langgraph\_graph.png)
    \item[\textcolor{green}{\checkmark}] Figure 2.6 : ReAct vs LangGraph — \textbf{À CRÉER} (tableau comparatif visuel)
    \item[\textcolor{green}{\checkmark}] Figure 2.7 : Architecture fonctionnelle — \textbf{EXISTANT} (fig\_03\_global\_architecture.png)
    \item[\textcolor{green}{\checkmark}] Figure 2.8 : Architecture logique — \textbf{EXISTANT} (fig\_03\_global\_architecture.png peut être dupliqué)
    \item[\textcolor{green}{\checkmark}] Figure 2.9 : Gantt — \textbf{EXISTANT} (fig\_01\_gantt.png)
    \item[\textcolor{green}{\checkmark}] Figure 2.10 : UC global — \textbf{EXISTANT} (fig\_03\_use\_case\_global.png)
    \item[\textcolor{green}{\checkmark}] Figure 2.11 : DoD — \textbf{À CRÉER} (diagramme ou tableau)
    \item[\textcolor{orange}{!}] Figures 2.12-2.20 : Logos technologies — \textbf{À CRÉER} (récupérer logos officiels + rédiger descriptions)
    
    \item[\textcolor{green}{\checkmark}] Tableau 2.1 : Comparatif modèles — \textbf{À CRÉER} (LLM/RAG/LinUCB/classique)
    \item[\textcolor{green}{\checkmark}] Tableau 2.2 : Product Backlog — \textbf{À CRÉER} (User Stories priorisées)
    \item[\textcolor{green}{\checkmark}] Tableau 2.3 : Planification Sprints — \textbf{EXISTANT} (dans chapitre\_02\_lancement.tex)
    \item[\textcolor{green}{\checkmark}] Tableau 2.4 : Ordinateur dev — \textbf{À CRÉER} (specs techniques)
\end{itemize}


% ============================================================================
% CHAPITRE 3 — SPRINT 1 : MISE EN PLACE ENVIRONNEMENT ET INTERFACE
% ============================================================================
\section{Chapitre 3 : Sprint 1 — Mise en place environnement et interface utilisateur}

\subsection*{Figures attendues}

\begin{longtable}{p{2cm} p{12cm}}
    \textbf{Référence} & \textbf{Description} \\
    \midrule
    \endhead
    
    Figure 3.1 & Diagramme de cas d'utilisation "Sprint 1" (authentification + health check) \\
    Figure 3.2 & Diagramme de séquence "Inscription/Connexion utilisateur" \\
    Figure 3.3 & Diagramme de séquence "Vérification état santé système" \\
    Figure 3.4 & Capture d'écran : Interface de connexion/inscription \\
    Figure 3.5 & Capture d'écran : Dashboard administrateur (après connexion) \\
    Figure 3.6 & Capture d'écran : Endpoint /health dans Swagger UI \\
    Figure 3.7 & Capture d'écran : Réponse JSON de /health \\
    Figure 3.8 & Capture d'écran : Schéma de base de données (pgAdmin) \\
    Figure 3.9 & Capture d'écran : Configuration environnement (Docker/Poetry) \\
    Figure 3.10 & Capture d'écran : Dépôt Git avec structure de projet \\
    \bottomrule
\end{longtable}

\subsection*{Tableaux attendus}

\begin{longtable}{p{2cm} p{12cm}}
    \textbf{Référence} & \textbf{Description} \\
    \midrule
    \endhead
    
    Tableau 3.1 & Sprint BackLog du Sprint 1 (ID User Story / ID Task / Task / Complexité XS/S/M/L/XL) \\
    Tableau 3.2 & Description Textuelle du cas d'utilisation "S'authentifier" (Titre, Acteur, Pré-condition, Scénario nominal, Post-condition) \\
    Tableau 3.3 & Description Textuelle du cas d'utilisation "Consulter état santé" \\
    \bottomrule
\end{longtable}

\subsection*{Statut de création}

\begin{itemize}
    \item[\textcolor{green}{\checkmark}] Figure 3.1 : UC Sprint 1 — \textbf{À CRÉER} (sous-ensemble de fig\_03\_use\_case\_global.png)
    \item[\textcolor{green}{\checkmark}] Figure 3.2 : Séquence authentification — \textbf{À CRÉER} (diagramme UML)
    \item[\textcolor{green}{\checkmark}] Figure 3.3 : Séquence health check — \textbf{EXISTANT} (fig\_03\_security\_flow.png peut être adapté)
    \item[\textcolor{orange}{!}] Figures 3.4-3.10 : Captures écran — \textbf{PARTIELLEMENT EXISTANT} (screen\_06\_health\_endpoint.png existe, autres à capturer)
    
    \item[\textcolor{green}{\checkmark}] Tableau 3.1 : Sprint Backlog 1 — \textbf{EXISTANT} (dans chapitre\_03\_sprint1.tex)
    \item[\textcolor{green}{\checkmark}] Tableau 3.2 : UC Authentification — \textbf{À CRÉER} (description textuelle)
    \item[\textcolor{green}{\checkmark}] Tableau 3.3 : UC Health Check — \textbf{EXISTANT} (tab:uc\_health dans chapitre\_03\_sprint1.tex)
\end{itemize}


% ============================================================================
% CHAPITRE 4 — SPRINT 2 : IMPLÉMENTATION DU PIPELINE RAG ET NLP
% ============================================================================
\section{Chapitre 4 : Sprint 2 — Pipeline RAG et traitement du langage}

\subsection*{Figures attendues}

\begin{longtable}{p{2cm} p{12cm}}
    \textbf{Référence} & \textbf{Description} \\
    \midrule
    \endhead
    
    Figure 4.1 & Organigramme du Sprint 2 (flux de traitement RAG) \\
    Figure 4.2 & Diagramme de séquence du Sprint 2 (requête → embedding → retrieval → reranking) \\
    Figure 4.3 & Diagramme de classes : Entités RAG (Document, Chunk, Embedding) \\
    Figure 4.4 & Capture d'écran : Interface de test du pipeline RAG \\
    Figure 4.5 & Capture d'écran : Résultats de recherche vectorielle (top-k chunks) \\
    Figure 4.6 & Capture d'écran : Scores de reranking (cross-encoder) \\
    Figure 4.7 & Capture d'écran : Logs de réindexation incrémentale \\
    Figure 4.8 & Capture d'écran : Interface admin de rebuild d'index \\
    \bottomrule
\end{longtable}

\subsection*{Tableaux attendus}

\begin{longtable}{p{2cm} p{12cm}}
    \textbf{Référence} & \textbf{Description} \\
    \midrule
    \endhead
    
    Tableau 4.1 & Sprint BackLog du Sprint 2 (ID User Story / ID Task / Task / Complexité) \\
    Tableau 4.2 & Description Textuelle du cas d'utilisation "Exécuter recherche RAG" \\
    Tableau 4.3 & Comparaison des modèles d'embedding (BAAI/bge-m3 vs autres) \\
    \bottomrule
\end{longtable}

\subsection*{Statut de création}

\begin{itemize}
    \item[\textcolor{green}{\checkmark}] Figure 4.1 : Organigramme RAG — \textbf{EXISTANT} (fig\_02\_rag\_architecture.png peut être adapté)
    \item[\textcolor{green}{\checkmark}] Figure 4.2 : Séquence RAG — \textbf{À CRÉER} (diagramme UML)
    \item[\textcolor{green}{\checkmark}] Figure 4.3 : Classes RAG — \textbf{À CRÉER} (diagramme UML)
    \item[\textcolor{orange}{!}] Figures 4.4-4.8 : Captures écran — \textbf{PARTIELLEMENT EXISTANT} (screen\_04\_rag\_search\_result.png existe, autres à capturer)
    
    \item[\textcolor{green}{\checkmark}] Tableau 4.1 : Sprint Backlog 2 — \textbf{À CRÉER} (dans chapitre\_04\_sprint2.tex)
    \item[\textcolor{green}{\checkmark}] Tableau 4.2 : UC Recherche RAG — \textbf{À CRÉER} (description textuelle)
    \item[\textcolor{green}{\checkmark}] Tableau 4.3 : Embeddings — \textbf{À CRÉER} (comparatif technique)
\end{itemize}


% ============================================================================
% CHAPITRE 5 — SPRINT 3 : AGENT LLM ET PLANIFICATION CONVERSATIONNELLE
% ============================================================================
\section{Chapitre 5 : Sprint 3 — Agent LLM et planification conversationnelle}

\subsection*{Figures attendues}

\begin{longtable}{p{2cm} p{12cm}}
    \textbf{Référence} & \textbf{Description} \\
    \midrule
    \endhead
    
    Figure 5.1 & Processus de l'agent LangGraph (équivalent "Processus de prédiction des intentions") \\
    Figure 5.2 & Diagramme de séquence du Sprint 3 "Agent planificateur" \\
    Figure 5.3 & Architecture détaillée des couches du modèle LangGraph (THINK/EXECUTE/VALIDATE/RESPOND) \\
    Figure 5.4 & Diagramme d'activité : Cycle de vie d'une requête utilisateur \\
    Figure 5.5 & Capture d'écran : Interface de chat avec l'agent \\
    Figure 5.6 & Capture d'écran : Plan de voyage généré (sortie agent) \\
    Figure 5.7 & Capture d'écran : Étapes de réflexion (node THINK) \\
    Figure 5.8 & Capture d'écran : Outils appelés par l'agent (logs) \\
    Figure 5.9 & Capture d'écran : Validation des sorties (node VALIDATE) \\
    \bottomrule
\end{longtable}

\subsection*{Tableaux attendus}

\begin{longtable}{p{2cm} p{12cm}}
    \textbf{Référence} & \textbf{Description} \\
    \midrule
    \endhead
    
    Tableau 5.1 & Sprint BackLog du Sprint 3 (ID User Story / ID Task / Task / Complexité) \\
    Tableau 5.2 & Description Textuelle du cas d'utilisation "Générer un plan de voyage" \\
    Tableau 5.3 & Liste des 18 outils LLM-callable (nom, description, paramètres) \\
    \bottomrule
\end{longtable}

\subsection*{Statut de création}

\begin{itemize}
    \item[\textcolor{green}{\checkmark}] Figure 5.1 : Processus LangGraph — \textbf{EXISTANT} (fig\_02\_langgraph\_graph.png)
    \item[\textcolor{green}{\checkmark}] Figure 5.2 : Séquence agent — \textbf{EXISTANT} (fig\_02\_agent\_sequence.png)
    \item[\textcolor{green}{\checkmark}] Figure 5.3 : Couches LangGraph — \textbf{À CRÉER} (détail des 4 nœuds)
    \item[\textcolor{green}{\checkmark}] Figure 5.4 : Activité requête — \textbf{EXISTANT} (fig\_04\_langgraph\_routing.png peut être adapté)
    \item[\textcolor{orange}{!}] Figures 5.5-5.9 : Captures écran — \textbf{PARTIELLEMENT EXISTANT} (screen\_04\_chat\_plan\_output.png, screen\_06\_plan\_thinking\_steps.png existent)
    
    \item[\textcolor{green}{\checkmark}] Tableau 5.1 : Sprint Backlog 3 — \textbf{À CRÉER} (dans chapitre\_05\_sprint3.tex)
    \item[\textcolor{green}{\checkmark}] Tableau 5.2 : UC Générer plan — \textbf{À CRÉER} (description textuelle)
    \item[\textcolor{green}{\checkmark}] Tableau 5.3 : 18 outils — \textbf{À CRÉER} (catalogue des tools)
\end{itemize}


% ============================================================================
% CHAPITRE 6 — SPRINT 4 : MOTEUR DE RECOMMANDATION LINUCB
% ============================================================================
\section{Chapitre 6 : Sprint 4 — Moteur de recommandation LinUCB}

\subsection*{Figures attendues}

\begin{longtable}{p{2cm} p{12cm}}
    \textbf{Référence} & \textbf{Description} \\
    \midrule
    \endhead
    
    Figure 6.1 & Processus du LinUCB (équivalent "Processus de l'inverted index") \\
    Figure 6.2 & Organigramme du Sprint 4 "Générer les recommandations" \\
    Figure 6.3 & Diagramme de cas d'utilisation "Générer les recommandations" \\
    Figure 6.4 & Diagramme de séquence "Générer recommandation" \\
    Figure 6.5 & Diagramme de séquence "Collecter événement comportemental" \\
    Figure 6.6 & Architecture du vecteur de caractéristiques (produit de Kronecker) \\
    Figure 6.7 & Pipeline de classement en 10 étapes \\
    Figure 6.8 & Capture d'écran : Interface de recommandations avec résultats \\
    Figure 6.9 & Capture d'écran : Événements comportementaux collectés (logs) \\
    Figure 6.10 & Capture d'écran : Matrices LinUCB en mémoire \\
    Figure 6.11 & Graphique : Cumulative reward sur période de test \\
    \bottomrule
\end{longtable}

\subsection*{Tableaux attendus}

\begin{longtable}{p{2cm} p{12cm}}
    \textbf{Référence} & \textbf{Description} \\
    \midrule
    \endhead
    
    Tableau 6.1 & Sprint BackLog du Sprint 4 (ID User Story / ID Task / Task / Complexité) \\
    Tableau 6.2 & Description Textuelle de cas d'utilisation "Générer les recommandations" \\
    Tableau 6.3 & Signaux comportementaux (8 types de récompenses) \\
    Tableau 6.4 & Performance de scoring (latence avant/après cache) \\
    \bottomrule
\end{longtable}

\subsection*{Statut de création}

\begin{itemize}
    \item[\textcolor{green}{\checkmark}] Figure 6.1 : Processus LinUCB — \textbf{EXISTANT} (fig\_02\_linucb\_cycle.png)
    \item[\textcolor{green}{\checkmark}] Figure 6.2 : Organigramme Sprint 4 — \textbf{À CRÉER} (flux de ranking)
    \item[\textcolor{green}{\checkmark}] Figure 6.3 : UC Recommandations — \textbf{À CRÉER} (diagramme UML)
    \item[\textcolor{green}{\checkmark}] Figure 6.4 : Séquence génération — \textbf{À CRÉER} (diagramme UML)
    \item[\textcolor{green}{\checkmark}] Figure 6.5 : Séquence collecte — \textbf{EXISTANT} (fig\_05\_tracker\_flush.png peut être adapté)
    \item[\textcolor{green}{\checkmark}] Figure 6.6 : Vecteur Kronecker — \textbf{À CRÉER} (schéma mathématique)
    \item[\textcolor{green}{\checkmark}] Figure 6.7 : Pipeline ranking — \textbf{EXISTANT} (fig\_05\_ranker\_pipeline.png)
    \item[\textcolor{orange}{!}] Figures 6.8-6.10 : Captures écran — \textbf{PARTIELLEMENT EXISTANT} (screen\_A4\_recommendations\_response.png existe)
    \item[\textcolor{green}{\checkmark}] Figure 6.11 : Cumulative reward — \textbf{EXISTANT} (fig\_06\_cumulative\_reward.png)
    
    \item[\textcolor{green}{\checkmark}] Tableau 6.1 : Sprint Backlog 4 — \textbf{À CRÉER} (dans chapitre\_06\_sprint4.tex)
    \item[\textcolor{green}{\checkmark}] Tableau 6.2 : UC Recommandations — \textbf{À CRÉER} (description textuelle)
    \item[\textcolor{green}{\checkmark}] Tableau 6.3 : Signaux comportementaux — \textbf{À CRÉER} (8 types de rewards)
    \item[\textcolor{green}{\checkmark}] Tableau 6.4 : Performance scoring — \textbf{À CRÉER} (benchmarks)
\end{itemize}


% ============================================================================
% CHAPITRE 7 — SPRINT 5 : TESTS, ÉVALUATION ET DÉPLOIEMENT
% ============================================================================
\section{Chapitre 7 : Sprint 5 — Tests, évaluation et déploiement}

\subsection*{Figures attendues}

\begin{longtable}{p{2cm} p{12cm}}
    \textbf{Référence} & \textbf{Description} \\
    \midrule
    \endhead
    
    Figure 7.1 & Diagramme de cas d'utilisation "Exécuter tests unitaires" \\
    Figure 7.2 & Diagramme de séquence "Déploiement CI/CD" \\
    Figure 7.3 & Organigramme de la stratégie de tests \\
    Figure 7.4 & Capture d'écran : Rapports de tests pytest (couverture) \\
    Figure 7.5 & Capture d'écran : Dashboard de monitoring (latences) \\
    Figure 7.6 & Capture d'écran : Interface de déploiement (Docker/Kubernetes) \\
    Figure 7.7 & Graphique : Latence moyenne par endpoint API \\
    Figure 7.8 & Graphique : Taux de satisfaction utilisateur (enquête) \\
    \bottomrule
\end{longtable}

\subsection*{Tableaux attendus}

\begin{longtable}{p{2cm} p{12cm}}
    \textbf{Référence} & \textbf{Description} \\
    \midrule
    \endhead
    
    Tableau 7.1 & Sprint BackLog du Sprint 5 (ID User Story / ID Task / Task / Complexité) \\
    Tableau 7.2 & Résultats des tests unitaires (nombre de tests, couverture, échecs) \\
    Tableau 7.3 & Métriques de performance (latence, throughput, disponibilité) \\
    Tableau 7.4 & Checklist de déploiement (prérequis, validation, rollback) \\
    \bottomrule
\end{longtable}

\subsection*{Statut de création}

\begin{itemize}
    \item[\textcolor{green}{\checkmark}] Figure 7.1 : UC Tests — \textbf{À CRÉER} (diagramme UML)
    \item[\textcolor{green}{\checkmark}] Figure 7.2 : Séquence CI/CD — \textbf{À CRÉER} (diagramme UML)
    \item[\textcolor{green}{\checkmark}] Figure 7.3 : Stratégie tests — \textbf{À CRÉER} (organigramme)
    \item[\textcolor{orange}{!}] Figures 7.4-7.8 : Captures/Graphiques — \textbf{À CAPTURER} (selon environnement de test)
    
    \item[\textcolor{green}{\checkmark}] Tableau 7.1 : Sprint Backlog 5 — \textbf{À CRÉER} (dans chapitre\_07\_sprint5.tex)
    \item[\textcolor{green}{\checkmark}] Tableau 7.2 : Résultats tests — \textbf{À CRÉER} (rapports pytest)
    \item[\textcolor{green}{\checkmark}] Tableau 7.3 : Métriques performance — \textbf{À CRÉER} (benchmarks finaux)
    \item[\textcolor{green}{\checkmark}] Tableau 7.4 : Checklist déploiement — \textbf{À CRÉER} (procédure)
\end{itemize}


% ============================================================================
% CONCLUSION GÉNÉRALE
% ============================================================================
\section{Conclusion Générale et Perspectives}

\subsection*{Éléments attendus}

\begin{itemize}
    \item Synthèse des travaux réalisés par Sprint
    \item Récapitulatif des objectifs atteints (vs Product Backlog initial)
    \item Limitations identifiées (techniques, temporelles, fonctionnelles)
    \item Perspectives d'amélioration (court terme / moyen terme / long terme)
    \item Ouverture vers des travaux futurs (research directions)
\end{itemize}


% ============================================================================
% ANNEXES
% ============================================================================
\section{Annexes}

\subsection*{Figures annexes}

\begin{longtable}{p{2cm} p{12cm}}
    \textbf{Référence} & \textbf{Description} \\
    \midrule
    \endhead
    
    Figure A.1 & Schéma complet de la base de données (MCD/MLD) \\
    Figure A.2 & Exemples de conversations avec l'agent (scénarios types) \\
    Figure A.3 & Configuration des modèles LLM (prompts système) \\
    Figure A.4 & Captures d'écran de l'interface complète (vue utilisateur) \\
    Figure A.5 & API documentation (Swagger/OpenAPI) \\
    \bottomrule
\end{longtable}

\subsection*{Tableaux annexes}

\begin{longtable}{p{2cm} p{12cm}}
    \textbf{Référence} & \textbf{Description} \\
    \midrule
    \endhead
    
    Tableau A.1 & Glossaire des termes techniques (LLM, RAG, LinUCB, etc.) \\
    Tableau A.2 & Liste des abréviations utilisées dans le rapport \\
    Tableau A.3 & Références bibliographiques complètes (format IEEE) \\
    Tableau A.4 & Code source des endpoints API principaux \\
    \bottomrule
\end{longtable}

\subsection*{Statut de création}

\begin{itemize}
    \item[\textcolor{green}{\checkmark}] Figure A.1 : Schéma DB — \textbf{EXISTANT} (fig\_03\_er\_diagram\_planner.png, fig\_03\_er\_diagram\_recommender.png)
    \item[\textcolor{green}{\checkmark}] Figure A.2 : Conversations — \textbf{À CRÉER} (extraits de logs)
    \item[\textcolor{green}{\checkmark}] Figure A.3 : Prompts — \textbf{À CRÉER} (system prompts)
    \item[\textcolor{orange}{!}] Figure A.4 : Interface complète — \textbf{PARTIELLEMENT EXISTANT} (screen\_A4\_*.png existent partiellement)
    \item[\textcolor{green}{\checkmark}] Figure A.5 : API docs — \textbf{EXISTANT} (screen\_A4\_api\_docs.png)
    
    \item[\textcolor{green}{\checkmark}] Tableau A.1 : Glossaire — \textbf{EXISTANT} (dans chapter\_00\_preliminaries.tex)
    \item[\textcolor{green}{\checkmark}] Tableau A.2 : Abréviations — \textbf{EXISTANT} (dans chapter\_00\_preliminaries.tex)
    \item[\textcolor{green}{\checkmark}] Tableau A.3 : Bibliographie — \textbf{EXISTANT} (references.bib)
    \item[\textcolor{green}{\checkmark}] Tableau A.4 : Code source — \textbf{À CRÉER} (extraits significatifs)
\end{itemize}


% ============================================================================
% RÉCAPITULATIF GLOBAL
% ============================================================================
\newpage
\section{Récapitulatif Global des Figures et Tableaux}

\subsection*{Total des figures par chapitre}

\begin{longtable}{p{4cm} r r r}
    \textbf{Chapitre} & \textbf{Figures existantes} & \textbf{Figures à créer} & \textbf{Total} \\
    \midrule
    \endhead
    
    Chapitre 1 (Contexte) & 5 & 6 & 11 \\
    Chapitre 2 (Lancement) & 9 & 11 & 20 \\
    Chapitre 3 (Sprint 1) & 2 & 8 & 10 \\
    Chapitre 4 (Sprint 2) & 2 & 6 & 8 \\
    Chapitre 5 (Sprint 3) & 5 & 4 & 9 \\
    Chapitre 6 (Sprint 4) & 5 & 6 & 11 \\
    Chapitre 7 (Sprint 5) & 0 & 8 & 8 \\
    Annexes & 3 & 2 & 5 \\
    \midrule
    \textbf{TOTAL} & \textbf{31} & \textbf{51} & \textbf{82} \\
    \bottomrule
\end{longtable}

\subsection*{Total des tableaux par chapitre}

\begin{longtable}{p{4cm} r r r}
    \textbf{Chapitre} & \textbf{Tableaux existants} & \textbf{Tableaux à créer} & \textbf{Total} \\
    \midrule
    \endhead
    
    Chapitre 1 (Contexte) & 0 & 3 & 3 \\
    Chapitre 2 (Lancement) & 1 & 3 & 4 \\
    Chapitre 3 (Sprint 1) & 2 & 1 & 3 \\
    Chapitre 4 (Sprint 2) & 0 & 3 & 3 \\
    Chapitre 5 (Sprint 3) & 0 & 3 & 3 \\
    Chapitre 6 (Sprint 4) & 0 & 4 & 4 \\
    Chapitre 7 (Sprint 5) & 0 & 4 & 4 \\
    Annexes & 3 & 1 & 4 \\
    \midrule
    \textbf{TOTAL} & \textbf{6} & \textbf{25} & \textbf{31} \\
    \bottomrule
\end{longtable}

\subsection*{Priorisation des créations}

\textbf{Priorité 1 (Critique - avant rédaction finale) :}
\begin{itemize}
    \item Toutes les figures de diagrammes UML (cas d'utilisation, séquence, activité)
    \item Tous les Sprint Backlog (tableaux avec complexité XS/S/M/L/XL)
    \item Definition of Done (Figure 2.11)
    \item Product Backlog complet (Tableau 2.2)
\end{itemize}

\textbf{Priorité 2 (Important - pendant rédaction) :}
\begin{itemize}
    \item Captures d'écran des interfaces principales
    \item Logos des technologies avec descriptions
    \item Tableaux comparatifs (méthodologies, modèles, embeddings)
\end{itemize}

\textbf{Priorité 3 (Secondaire - relecture finale) :}
\begin{itemize}
    \item Graphiques de performance et métriques
    \item Photos de l'équipe projet
    \item Captures d'écran secondaires (logs, configuration)
\end{itemize}


% ============================================================================
% CONCLUSION DU DOCUMENT
% ============================================================================
\section*{Conclusion}

Ce document présente la liste exhaustive des 82 figures et 31 tableaux requis pour le rapport de PFE Guidni, structurés exactement selon le modèle du rapport d'Amal Rahmeni. 

\textbf{État actuel :} 31 figures et 6 tableaux sont déjà existants ou partiellement existants dans le dépôt.

\textbf{Travail restant :} 51 figures et 25 tableaux doivent être créés, principalement :
\begin{itemize}
    \item Diagrammes UML complémentaires (séquence, activité, cas d'utilisation par Sprint)
    \item Tableaux de Sprint Backlog pour chaque itération
    \item Captures d'écran des interfaces utilisateur
    \item Logos des technologies avec descriptions académiques
    \item Graphiques de performance et métriques d'évaluation
\end{itemize}

La numérotation suit strictement le format X.Y où X = numéro du chapitre et Y = numéro dans le chapitre, conformément aux standards académiques français.

\end{document}
